% =============================================================================
% Kapitel 1: Einleitung
% =============================================================================
\chapter{Einleitung}
\label{ch:einleitung}

Die gezielte Entwicklung neuer Funktionsmaterialien setzt ein genaues
Verständnis ihrer atomaren Struktur voraus.
Während Beugungsmethoden wie Röntgendiffraktometrie die langreichweitige
Kristallstruktur beschreiben, fehlt ihnen in der Regel die Empfindlichkeit
für lokale Abweichungen von der mittleren Ordnung.
Gerade bei Mehrkomponentensystemen -- Hochentropie- und
Mittelentropielegierungen, bimetallischen Nanopartikeln oder dotierten
Oxiden -- bestimmen jedoch diese lokalen Arrangements wesentlich die
mechanischen, katalytischen und magnetischen Eigenschaften
\cite{Yeh2004, George2019, Ferrando2008}.

Die Röntgenabsorptionsfeinstrukturspektroskopie (EXAFS) bietet hier einen
komplementären Zugang.
Sie nutzt die Modulation des Absorptionskoeffizienten oberhalb einer
elementspezifischen Absorptionskante, um Bindungslängen,
Koordinationszahlen und die thermische bzw.\ statische Unordnung in der
unmittelbaren Nachbarschaft eines Absorberatoms zu bestimmen
\cite{Stern1974, Rehr2000}.
Da jede Kante einzeln adressiert wird, lässt sich die lokale Umgebung
jedes Elements in einem Mehrstoffsystem separat untersuchen.

\section{Motivation}

Bei einer ternären Legierung wie CoNiCu existieren neun verschiedene
Nachbarpaare (Co-Co, Co-Ni, Co-Cu, Ni-Co,\,\ldots).
Jeder dieser Beiträge wird durch mindestens vier Parameter beschrieben
(Koordinationszahl, Abstand, Debye-Waller-Faktor, Energieverschiebung),
sodass bereits für die erste Koordinationsschale $9 \times 4 = 36$
Fitparameter anfallen.
Die Menge an unabhängiger Information, die ein EXAFS-Spektrum enthält,
ist jedoch durch die sogenannte Nyquist-Zahl
\begin{equation}
  N_{\text{ind}} = \frac{2\,\Delta k\,\Delta R}{\pi} + 2
  \label{eq:nyquist}
\end{equation}
begrenzt, wobei $\Delta k$ den nutzbaren $k$-Bereich und $\Delta R$ das
betrachtete $R$-Fenster bezeichnen \cite{Stern1993}.
Für typische Werte ($\Delta k \approx \SI{8}{\per\Angstrom}$,
$\Delta R \approx \SI{2}{\Angstrom}$) ergibt sich
$N_{\text{ind}} \approx 12$ pro Kante, also insgesamt rund 36 unabhängige
Datenpunkte für drei Kanten -- kaum mehr als die Zahl der Parameter.
Ein konventioneller Least-Squares-Fit wird unter diesen Bedingungen
numerisch instabil und physikalisch mehrdeutig.

\section{Der RMC/EA-Ansatz}

Eine alternative Strategie kehrt das Fitproblem um:
Statt analytische Funktionen an das Spektrum anzupassen, wird ein
atomistisches Strukturmodell konstruiert und solange verändert, bis die
daraus berechneten Spektren aller Kanten gleichzeitig mit dem Experiment
übereinstimmen.
Die Software \evax{} implementiert diesen Ansatz als Kombination eines
evolutionären Algorithmus (EA) mit Reverse Monte Carlo (RMC)
\cite{Timoshenko2009, Timoshenko2012a, Timoshenko2014}.
Das Verfahren kommt ohne explizite Parametrisierung der
Paarverteilungsfunktionen aus und liefert unmittelbar
dreidimensionale Strukturmodelle, aus denen sich Koordinationszahlen,
Bindungslängen und die chemische Nahordnung ableiten lassen.

\section{Zielsetzung}

Ziel dieser Arbeit ist es, \evax{} auf zwei experimentelle Datensätze
anzuwenden:

\begin{enumerate}[label=(\roman*)]
  \item Eine \textbf{\TODO{CoNiCu}-Mittelentropielegierung (MEA)}, gemessen
        an den K-Kanten von Cobalt, Nickel und Kupfer.
        Durch gleichzeitige Auswertung aller drei Kanten sollen die
        elementaufgelösten Koordinationszahlen und insbesondere die
        Warren-Cowley-Parameter der chemischen Nahordnung bestimmt werden.
  \item \textbf{\TODO{CoPt}-Nanopartikel}, bei denen die Frage im
        Vordergrund steht, ob die beiden Elemente statistisch auf
        einem gemeinsamen Gitter verteilt sind oder eine geordnete
        Phase (etwa L1$_0$) bilden.
\end{enumerate}

Darüber hinaus wird die Sensitivität der Ergebnisse gegenüber den
Fit-Parametern ($k$-Bereich, $k$-Gewichtung, Vergleichsraum) durch
eine systematische Parameterstudie untersucht, und die Wahl der
Kristallstruktur (FCC vs.\ BCC vs.\ HCP) wird durch Vergleich der
Fitresiduen validiert.

\section{Aufbau der Arbeit}

Kapitel~\ref{ch:theorie} legt die physikalischen Grundlagen der
EXAFS-Spektroskopie dar.
In Kapitel~\ref{ch:evax} wird die Software \evax{} vorgestellt und ihr
Arbeitsablauf erläutert.
Kapitel~\ref{ch:methodik} beschreibt die konkrete Vorgehensweise,
einschließlich der Datenvorverarbeitung und der gewählten Fit-Strategie.
Die Ergebnisse der Auswertung beider Proben werden in
Kapitel~\ref{ch:auswertung} zusammengefasst und in
Kapitel~\ref{ch:diskussion} diskutiert.
Kapitel~\ref{ch:fazit} schließt die Arbeit mit einem Fazit und Ausblick ab.
